\documentclass[11pt]{article}
\usepackage[margin=1in]{geometry}
\usepackage[T1]{fontenc}
\usepackage[utf8]{inputenc}
\usepackage{lmodern}
\usepackage{microtype}
\usepackage{booktabs}
\usepackage{array}
\usepackage{enumitem}
\usepackage{hyperref}
\usepackage{natbib}

\title{Post-Training Amplifies Judge Framing Sensitivity: Draft Structure and Evidence Ledger}
\author{AISafety project draft}
\date{\today}

\begin{document}
\maketitle

\begin{abstract}
Large language model (LLM) judges often favor LLM-generated or LLM-presented
content over human-written alternatives. Existing work establishes this
phenomenon and related judge biases, but leaves open how much of the effect is
already present in base models and how much is introduced or amplified by
post-training for assistant-style use. This draft note organizes a workshop
paper around a narrower claim: base models already show substantial
LLM-favoring behavior, but instruction/post-training amplifies and reshapes
that behavior in binary judge settings by making models more sensitive to
evaluation framing and answer-like or formal presentation cues. The strongest
current evidence comes from a Tulu staged SFT/DPO ladder, Qwen 2.5 base-vs-IT
replication, deterministic surface-cue counterfactuals, matched
human-vs-LLM subsets, and likelihood-only negative controls. The remaining
claim gates are uncertainty estimates and one more cross-family/scale
replication, preferably Gemma 2 27B base/IT.
\end{abstract}

\section{Working Claim}

\paragraph{Main claim.}
Post-training does not create AI--AI or LLM-over-human preference from
scratch. In our experiments, base models already prefer LLM-authored text in
many paired settings. The post-training effect is better described as an
amplification and conditionalization of this pre-existing preference:
assistant-tuned models become more sensitive to binary judge framing and to
surface presentations that resemble helpful, formal, answer-like assistant
outputs.

\paragraph{Claim boundary.}
The current evidence supports a diagnostic account of judge behavior, not a
complete causal account of all AI--AI bias. The paper should avoid claiming
that RLHF alone explains LLM preference. It should instead claim that
post-training changes how model judges respond to evaluation prompts and
presentation cues, especially in forced-choice settings. Likelihood-only
controls are crucial because they show that the main forced-choice effects are
not simply lower perplexity or generic next-token familiarity.

\section{Intended Contributions}

\begin{enumerate}[leftmargin=1.5em]
\item \textbf{Stage contrast.} Compare base and post-trained checkpoints on
the same human-vs-LLM paired corpus, including a staged Tulu/Llama ladder
(base, SFT, DPO, final, Llama instruct) and Qwen 2.5 base/IT replication.

\item \textbf{Template interaction.} Show that post-training changes sensitivity
to judge prompt framing. Minimal prompts often reduce or collapse effects that
standard, rubric, or substance-only judge prompts recover.

\item \textbf{Deterministic surface-cue control.} Reuse deterministic
counterfactual pairings to test whether the same stage/template pattern appears
for controlled presentation cues rather than direct human-vs-LLM labels.

\item \textbf{Likelihood negative control.} Compare forced-choice margins with
response-likelihood scoring. The likelihood-only deltas are consistently small
or negative relative to forced-choice deltas.

\item \textbf{Confound controls.} Use matched pair subsets controlling length,
prompt overlap, type-token ratio, and punctuation, plus order swaps to reduce
first-item bias.

\item \textbf{Next claim gate.} Add paired bootstrap confidence intervals and a
Gemma 2 27B base/IT family/scale replication before submission.
\end{enumerate}

\section{Current Evidence Ledger}

\subsection{Human-vs-LLM Stage Contrast}

The broad and matched Tulu/Llama results show that base models already favor
LLM text, but Tulu post-training substantially increases LLM-minus-human
forced-choice margins. On the relaxed matched set of 1,370 pairs:

\begin{center}
\small
\begin{tabular}{lrrr}
\toprule
Contrast & Mean margin delta & Preference-rate movement & Interpretation \\
\midrule
Tulu SFT $-$ Llama base & $+0.747$ & $0.788 \to 0.885$ & Large SFT amplification \\
Tulu DPO $-$ Tulu SFT & $+0.620$ & $0.885 \to 0.874$ & Margin increase, fewer hard flips \\
Tulu final $-$ Llama base & $+1.279$ & $0.788 \to 0.880$ & Large full-stage effect \\
Llama instruct $-$ Llama base & $+0.661$ & $0.788 \to 0.758$ & Margin increase without rate increase \\
\bottomrule
\end{tabular}
\end{center}

Likelihood-only controls on the same matched set are much smaller:
Tulu SFT-like $-$ base-like is $+0.087$, and Tulu DPO-like $-$ SFT-like is
$+0.178$. This suggests that the large forced-choice effect is not merely
explained by the model assigning higher likelihood to LLM responses.

\subsection{Template Sensitivity on Human-vs-LLM Pairs}

Prompt templates tested so far are \texttt{minimal}, \texttt{standard},
\texttt{rubric\_quality}, and \texttt{substance\_only}. For Tulu on the matched
set, minimal prompts greatly reduce the base-to-SFT effect, while standard,
rubric, and substance prompts recover large effects. The key interaction is
that post-training itself increases template sensitivity. For
\texttt{standard} minus \texttt{minimal}, the Tulu SFT-minus-base interaction
is about $+0.756$, and Tulu final-minus-base is about $+1.005$.

Qwen 2.5 7B replicates the human-vs-LLM template result. With chat-template
prompting, Qwen Instruct shows a very large standard-minus-minimal interaction
relative to Qwen base ($+2.877$). When Qwen Instruct is forced through plain
prompting, the interaction remains large ($+1.407$), so the human-vs-LLM
template effect is amplified by chat formatting but not caused by it.

\subsection{Deterministic Surface-BT Counterfactuals}

The surface-BT suite converts deterministic cue-increase/decrease
counterfactuals into balanced binary comparisons. Axes tested:
\texttt{answer\_likeness\_packaging},
\texttt{formal\_institutional\_packaging}, and
\texttt{structured\_assistant\_packaging}. The current paired set contains
6,381 counterfactuals and 12,762 order-swapped BT rows.

For Tulu, surface-BT template sensitivity is positive and selective. It is
strongest for answer-likeness and formal/institutional packaging, holds for
both human-origin and LLM-origin texts, and is near zero under likelihood-only
scoring. This supports the interpretation that post-training increases
frame-conditioned sensitivity to assistant-like presentation cues rather than
only to the human/LLM source distinction.

For Qwen 2.5, the chat-template surface-BT run initially shows a strong
standard-minus-minimal interaction ($+0.676$), with axis-level interactions of
$+0.778$ for answer-likeness and $+0.692$ for formal/institutional packaging.
However, forcing Qwen Instruct to plain prompting removes the surface-BT
template interaction ($-0.023$). This is an important qualification:
Qwen's surface-cue template sensitivity is largely chat-template mediated,
although Qwen Instruct still shows a plain-prompt cue-plus margin increase over
base in both minimal ($+0.207$) and standard ($+0.184$) settings.

\section{Proposed Paper Structure}

\subsection{1. Introduction}

State the empirical problem: LLM judges are used widely, but they can favor
LLM-authored or LLM-presented content. Prior work establishes the phenomenon;
this paper asks what post-training contributes. Present the refined thesis:
base models already exhibit LLM preference, while post-training amplifies and
reshapes that preference in binary judge settings by increasing sensitivity to
judge framing and surface presentation cues.

\subsection{2. Related Work}

Organize related work around four clusters:

\paragraph{LLM-as-a-judge and known biases.}
MT-Bench and Chatbot Arena established LLM-as-a-judge as a scalable evaluation
approach and documented position, verbosity, and self-enhancement biases
\citep{zheng2023mtbench}. Position bias has since been studied more
systematically across judges and tasks \citep{shi2024positionbias}.

\paragraph{AI--AI bias and self-preference.}
AI--AI bias work shows that LLM systems can prefer LLM-generated
communications over human-authored alternatives in binary choice scenarios
\citep{laurito2025aiaibias}. Self-preference work shows that LLM judges can
favor their own or familiar generations and relates part of this to perplexity
or familiarity \citep{panickssery2024selfrecognition,
wataoka2024selfpreference}. Reward-model detection work similarly suggests
that aligned-model generations can be recognized by evaluators
\citep{lee2024remodetect}.

\paragraph{Preference leakage and model relatedness.}
Preference leakage studies bias arising when data generators and judges are
related by model identity, inheritance, or family membership
\citep{li2025preferenceleakage}. Our study is adjacent but distinct: we compare
base and post-trained checkpoints and ask how post-training changes judge
framing sensitivity, not only same-family preference.

\paragraph{Style and surface-presentation bias.}
Recent work emphasizes that LLM judges can be sensitive to style, formatting,
and epistemic markers rather than solely content quality
\citep{lee2024ember,soumik2026judging}. Our deterministic surface-BT suite
connects this line of work to base-vs-instruct stage contrasts and
likelihood-only controls.

\subsection{3. Data and Controls}

Describe HC3, HAP-E, H-LLMC2, and derived matched subsets. Explain source
stratification, order swaps, length/lexical matching, and why scalar reward
heads are avoided in the present paper except as background.

\subsection{4. Methods}

\paragraph{Forced-choice scoring.}
Define LLM-minus-human margin from answer-token logit pooling or normalized
binary choice probabilities. Emphasize that all results are pairwise and
order-averaged.

\paragraph{Likelihood-only scoring.}
Define response-likelihood margin. This is the negative control for a
perplexity/familiarity-only explanation.

\paragraph{Template variants.}
Define \texttt{minimal}, \texttt{standard}, \texttt{rubric\_quality}, and
\texttt{substance\_only}. Treat \texttt{standard} minus \texttt{minimal} as the
main test of judge-framing sensitivity.

\paragraph{Surface-BT counterfactuals.}
Define cue-plus/cue-minus pairs and the axes. Explain that these deterministic
counterfactuals test presentation sensitivity beyond human-vs-LLM authorship.

\subsection{5. Results}

Recommended result order:

\begin{enumerate}[leftmargin=1.5em]
\item Base models already prefer LLM responses.
\item Tulu SFT/DPO stages amplify forced-choice LLM margins.
\item Likelihood-only deltas are much smaller.
\item Template sensitivity is post-training dependent.
\item Surface-BT shows selective sensitivity to answer-like/formal cues.
\item Qwen replicates human-vs-LLM template sensitivity, but reveals that
surface-BT template sensitivity can be chat-template mediated.
\item Bootstrap CIs and Gemma 27B replication, once available.
\end{enumerate}

\subsection{6. Discussion}

Discuss the mechanism as an interaction between pretraining priors,
assistant-style post-training, and evaluation-prompt affordances. The central
interpretation should be: post-training makes models better binary evaluators
of prompts written like evaluations, but this also makes them more responsive
to surface signals that correlate with assistant outputs.

\subsection{7. Limitations}

\begin{itemize}[leftmargin=1.5em]
\item Base preference remains substantially unexplained.
\item Tulu gives staged SFT/DPO evidence, but Qwen and Gemma provide only
base/IT contrasts.
\item Surface counterfactuals are deterministic and may not cover natural
style variation.
\item Qwen surface-BT template sensitivity is chat-template mediated.
\item No human reannotation yet verifies that the cue-plus variants are equal
quality under human judgments.
\item Bootstrap uncertainty and one more family/scale replication are still
needed before submission.
\end{itemize}

\section{Planned Tables and Figures}

\begin{enumerate}[leftmargin=1.5em]
\item \textbf{Table 1: Human-vs-LLM stage contrast.} Rows: base, SFT, DPO,
final, instruct; columns: margin, preference rate, likelihood control.
\item \textbf{Table 2: Template interactions.} Rows: model family/stage
contrast; columns: minimal, standard, standard-minus-minimal interaction.
\item \textbf{Table 3: Surface-BT axes.} Rows: answer-likeness, formal,
structured; columns: stage delta and template interaction by family.
\item \textbf{Figure 1: Conceptual pipeline.} Base/IT model pair, paired
human-vs-LLM data, template variants, forced-choice vs likelihood controls.
\item \textbf{Figure 2: Interaction plot.} Template delta by stage for Tulu,
Qwen, and Gemma.
\item \textbf{Figure 3: Surface axis decomposition.} Interaction split by
axis and source role.
\end{enumerate}

\section{Remaining Experiments Before Submission}

\paragraph{Paired bootstrap.}
Add paired bootstrap confidence intervals over pair IDs for human-vs-LLM runs
and over counterfactual IDs for surface-BT runs. Report 95\% intervals for
stage deltas and stage-by-template interactions.

\paragraph{Gemma 2 27B base/IT.}
Run the minimal decisive matrix first: HLLM minimal/standard with likelihood
control, and Surface-BT minimal/standard with likelihood control. If the
chat-template run is positive, run the plain-prompt instruct control.

\paragraph{Optional size check.}
If Gemma 27B is expensive or ambiguous, add Gemma 2 9B or Qwen 2.5 3B as a
same-family size ablation. This is secondary to the 27B cross-family check.

\paragraph{Human spot check.}
For a small subset of surface-BT cue-plus/cue-minus pairs, ask human reviewers
or a blind annotation protocol whether cue-plus is actually better. This is not
required for the diagnostic claim, but would strengthen the interpretation
that the model responds to presentation rather than genuine quality.

\section{Draft Abstract Candidate}

LLM-based judges are increasingly used to evaluate open-ended model outputs,
but they can favor LLM-generated content over human-written alternatives. We
study whether this preference is already present in base language models or is
amplified by assistant-style post-training. Across matched human-vs-LLM pairs,
a staged Llama/Tulu ladder, Qwen 2.5 base/IT replication, deterministic
surface-cue counterfactuals, and likelihood-only controls, we find that base
models already show substantial LLM preference, but post-training substantially
increases forced-choice LLM margins and sensitivity to judge prompt framing.
The effect is strongest under standard evaluation templates and for
answer-like or formal presentation cues, while likelihood-only scoring shows
much smaller shifts. Qwen replication confirms human-vs-LLM template
sensitivity even under plain prompting, but also shows that some surface-cue
effects are mediated by chat-template formatting. These results suggest that
assistant post-training reshapes models into more framing-sensitive binary
evaluators, amplifying pre-existing LLM/style preferences in common
LLM-as-a-judge settings.

\bibliographystyle{plainnat}
\bibliography{references}

\end{document}
